\subsection*{Failure decomposition by trap}

Each evaluation was built around an analytical trap and has at least one inference fallacy the failing runs committed. The two describe the same design principle from different perspectives. A trap is indexed by data modality and names the processing or analysis choice on which an agent can go wrong. A fallacy is indexed by the form of the inference and names the faulty step from evidence to conclusion that the wrong choice produces. During authoring, every candidate evaluation was checked against both catalogs. A single evaluation therefore carries on average 2.6 trap categories and one or two fallacies. The 91 evaluations together covered 21 trap categories and captured 27 distinct fallacies.

Among traps (Figure~\ref{fig:traps}A), quality control, controls and artifacts were the hardest (10.6\%, 15 evaluations), followed by normalization and representation (12.6\%, 8) and threshold or resolution choices (16.5\%, 17). The three most common traps fell near the mean over models of 20.1\%: analyzing the wrong subset or compartment (20.6\%, 33), choosing the wrong replicate or unit of analysis (21.6\%, 39) and computing a statistic that does not match the claim (22.2\%, 32).

The hardest fallacies to avoid were a proxy read as the construct it stands for (4.5\%, 3 evaluations), denominator neglect (7.9\%, 3), equivocation (13.0\%, 3), the continuum fallacy (13.4\%, 5), scale or level confusion (13.7\%, 7) and confounding or an omitted variable (14.8\%, 20). Confounding is also the most common fallacy in the benchmark, present in more than one evaluation in five. Correlation read as causation (16.6\%) and argument from ignorance (17.6\%) followed, whereas the law of small numbers (20.0\%), Simpson's paradox (24.5\%), the significance fallacy (25.8\%) and reification (26.0\%) were intermediate, and the easiest fallacies were those of comparison and aggregation: a faulty baseline (33.9\%), an inconsistent comparison (36.5\%, 3 evaluations) and the ecological fallacy (37.1\%, 4). Grouping the rows into statistical and reasoning classes gave similar overall rates (22.4\% vs. 19.4\%), so the difficulty lies in specific inference forms rather than in one class.
